\documentclass[11pt]{article}

% ====================================================
% PACKAGES
% ====================================================
\usepackage[margin=1in, headheight=13.6pt]{geometry}
\usepackage[T1]{fontenc}
\usepackage{lmodern}
\usepackage{microtype}
\usepackage{setspace}
\usepackage{tcolorbox}
\usepackage{amsmath, amssymb, mathtools}
\usepackage{siunitx}

\usepackage{graphicx}
\usepackage{float}
\usepackage{subcaption}

\usepackage{enumitem}
\usepackage{booktabs}
\usepackage{longtable}
\usepackage{array}
\usepackage{multirow}

\usepackage{xcolor}
\usepackage{hyperref}

\usepackage{fancyhdr}
\usepackage{lastpage}

\usepackage{listings}
\usepackage{indentfirst}
\usepackage[justification=centering]{caption}

\pagestyle{fancy}
\fancyhf{}
\lhead{\Assignment\ |\ \course}
\rhead{\Name\ |\ \today}
\cfoot{Page \thepage\ of~\pageref{LastPage}}

\hypersetup{
    colorlinks,
    linkcolor={red!50!black},
    citecolor={blue!50!black},
    urlcolor={blue!80!black}
}

\newcommand{\Name}{Arael Anaya}
\newcommand{\Assignment}{2.1 Designing a Data Product for Your Robotic Mission}
\newcommand{\course}{Space Robotics}

\begin{document}

% ====================================================
% SECTION I: DATA TO BE COLLECTED
% ====================================================
\section*{I. Data to Be Collected}

\subsection*{a. Goal}

\noindent\textbf{Mission Objective:} Autonomously characterize Europa's ice shell thickness
along a surface traverse and identify an optimal cryobot insertion site without real-time
Earth support (REQ.03.01: ground-penetrating radar sounder resolves ice thickness to
$\pm$10\% for depths of 1--30~km at $\geq$100~m along-track spatial resolution).

\noindent\textbf{Data Product Goal:} Provide engineers and scientists with a time-aligned
record of radar sounder returns, rover localization, terrain context, and system health to
verify that the subsurface map meets REQ.03.01 and that the onboard site-selection
algorithm correctly identifies a candidate cryobot insertion point.

\noindent\textbf{Success Criterion:} Full traverse track covered at $\leq$100~m along-track
spacing, ice thickness accuracy within $\pm$10\% against calibration targets, and at least
one candidate insertion site flagged by the onboard algorithm (ice thickness $<$5~km,
surface slope $<$20$^\circ$, terrain roughness $<$0.7).


\subsection*{b. Data Collected by the Robot}

\begin{itemize}[leftmargin=1.5em]

    \item \textbf{Category 1 --- Radar Sounder (Primary Science):}
    \begin{itemize}
        \item Two-way travel time to basal reflector (ns)
        \item Received signal amplitude at basal reflector (dB)
        \item Estimated bulk ice dielectric permittivity (dimensionless)
        \item Derived ice shell thickness (m)
    \end{itemize}

    \item \textbf{Category 2 --- Rover Localization:}
    \begin{itemize}
        \item Position estimate from visual odometry and inter-agent ranging (m, local ENU)
        \item 1$\sigma$ position uncertainty (m)
        \item Rover heading (deg)
        \item Localization error flag (Boolean: uncertainty exceeds $\pm$0.5~m per REQ.02.01)
    \end{itemize}

    \item \textbf{Category 3 --- Terrain Context:}
    \begin{itemize}
        \item Surface slope at measurement site (deg)
        \item Terrain roughness index (dimensionless, 0--1)
        \item Surface temperature (K)
    \end{itemize}

    \item \textbf{Category 4 --- System Health:}
    \begin{itemize}
        \item Total ionizing dose accumulated by radar electronics (krad)
        \item Radar noise floor (dB)
        \item Inter-agent link packet loss rate (\%)
        \item Per-trace data quality flag (NOMINAL / CLUTTER / NOISE\_FLOOR / INVALID)
    \end{itemize}

    \item \textbf{Daily Summary:}
    \begin{itemize}
        \item Total traverse distance completed in sol (m)
        \item Valid radar trace count (count)
        \item Minimum observed ice thickness in sol (m)
        \item Candidate insertion sites flagged by onboard algorithm (count)
        \item Mean inter-agent consensus convergence time (min)
    \end{itemize}

\end{itemize}


% ====================================================
% SECTION II: TYPES OF MEASUREMENTS ON THE SPREADSHEET
% ====================================================
\section*{II. Types of Measurements on the Spreadsheet}

\noindent The spreadsheet is organized into multiple sheets to keep the dataset clear and usable.
Below is a breakdown of each sheet:

\begin{enumerate}

    \item \textbf{Sheet 1 --- Radar\_Traces:}\\
    Time-aligned record of all sounder measurements, one row per trace. Records the
    along-track position, two-way travel time, received amplitude, estimated permittivity,
    and derived ice thickness (computed as $d = c \cdot t / 2\sqrt{\varepsilon_r}$), alongside
    a per-trace quality flag. This is the primary science record that feeds the
    subsurface map and site-selection logic.

    \item \textbf{Sheet 2 --- Localization\_Log:}\\
    Per-trace rover position fix synchronized to the Radar\_Traces timestamp index.
    Captures position (X, Y), heading, 1$\sigma$ uncertainty, and the REQ.02.01 error flag.
    Georeferences each sounder trace to a physical surface coordinate.

    \item \textbf{Sheet 3 --- Terrain\_and\_Health:}\\
    Per-trace terrain and system health data, aligned to the same trace index. Captures
    surface slope, roughness, temperature, accumulated TID, radar noise floor, and
    inter-agent packet loss. Contextualizes each radar measurement and flags traces
    collected under degraded system or terrain conditions.

    \item \textbf{Sheet 4 --- Daily Summary:}\\
    One row per sol. Auto-populated via formulas from Sheets 1--3: total traverse
    distance, valid trace count, minimum ice thickness, candidate site count, and mean
    consensus convergence time. Primary operational report for mission planning and
    onboard replan decisions.

\end{enumerate}


% ====================================================
% SECTION IV: RATIONALE & ALIGNMENT WITH MISSION
% ====================================================
\section*{IV. Rationale and Alignment with Mission Objective}

\noindent\textbf{Alignment with Objective:}
Each data category maps to a specific requirement. The radar sounder returns (Category 1)
are the direct measurement required by REQ.03.01. Without travel time and permittivity
there is no thickness inversion and no subsurface map. The localization log (Category 2)
is required by REQ.02.01: each trace must be placed within $\pm$0.5~m to produce a
spatially accurate map rather than an unordered list of depth readings. Terrain context
(Category 3) drives the site-selection filter, where slope and roughness constrain physical
accessibility independent of ice thickness. System health data (Category 4) is required
by REQ.01.02 and REQ.04.01: elevated TID can degrade radar performance, and packet
loss above 30\% directly stresses the consensus convergence bound.

\noindent\textbf{Why These Measurements Are Sufficient:}
The four categories together close the loop on REQ.03.01 and its dependencies. Travel
time plus permittivity produces thickness; localization georeferences it; terrain context
filters by physical accessibility; health flags bound measurement confidence. The daily
summary aggregates these into the per-sol statistics the onboard planner needs to decide
whether to continue traversing or commit to a cryobot insertion attempt.

\noindent\textbf{Potential Gaps or Limitations:}
The main limitation is permittivity estimation uncertainty when brine inclusions or
heterogeneous layering are present (RISK-05 in the requirements document). Brine lenses
increase local permittivity and produce reflections that can cause the thickness
algorithm to underestimate depth to the basal interface. The CLUTTER quality flag
mitigates this by excluding affected traces from site selection, but does not resolve the
ambiguity without a secondary instrument. Additionally, the along-track resolution of
100~m defined by REQ.03.01 produces a 1D profile along the traverse; cross-track
coverage depends on orbital relay scan width, which is not yet specified.


% ====================================================
% AI DISCLOSURE
% ====================================================
\vspace{1em}
\noindent\small\textit{AI Disclosure: Artificial intelligence tools were used to verify that the technical 
statements, parameter values, and requirement references made in this document are consistent with the background 
research on Europa ice-shell radar sounding, cryobot mission concepts, and planetary robotics requirements 
engineering. AI also generated fake data to demonstrate the data being recorded.}

\end{document}